\clearpage\thispagestyle{empty}
\pdfbookmark[0]{Quellen zu Gauß’ Brief an Bode vom September 1802}{bode18020910-sources}
\begin{center}\textbf{Quellen zu Gauß’ Brief an Bode vom September 1802}\end{center}
\noindent Diese KI-Übertragung in heutiges Deutsch folgt dem Brief Carl Friedrich Gauß’ an Johann Elert Bode vom 10. September 1802 in \href{https://gdz.sub.uni-goettingen.de/id/PPN236007467}{Gauß’ \textit{Werken}, Band~VI} (Göttingen, 1874), S.~223.
\par\smallskip\noindent Der \href{https://api.digitale-sammlungen.de/iiif/image/v2/bsb10538334_00237/full/1450,/0/default.jpg}{Erstdruck} steht auf S.~227--229 in Bodes \href{https://wwwadd.zah.uni-heidelberg.de/publikationen/baj/baj1805/baj1805.htm}{\textit{Astronomischem Jahrbuch für das Jahr 1805}} (Berlin, 1802). Der Brief gibt die Elemente für die Pallas-Ephemeride an, erörtert die Lichtstärke des Planeten und enthält eine Anmerkung zu den Vorhersagen nach den Elementen III. Die \href{https://api.digitale-sammlungen.de/iiif/image/v2/bsb10538334_00238/full/1450,/0/default.jpg}{Fußnote auf S.~228} ist im Jahrbuch mit G. unterzeichnet.
\par\smallskip\noindent Das Jahrbuch druckt die Ephemeride mit den \href{https://api.digitale-sammlungen.de/iiif/image/v2/bsb10538334_00239/full/1450,/0/default.jpg}{Oppositionszeiten auf S.~229}. In den \textit{Werken} steht eine Beschreibung der Ephemeride in eckigen Klammern. Die beiden Zeilen zur Opposition sind dort abgedruckt. Beide Übertragungen folgen dem Umfang der \textit{Werke} und bewahren die Zahlen der Neigungszeile ohne Einheitenzeichen.
\par\smallskip\noindent Übertragung und Satz wurden mit KI erstellt. Die Digitalisate des Jahrbuchs stammen von der Bayerischen Staatsbibliothek.
\clearpage
